\documentclass[10pt,a4paper]{amsart}
\usepackage[margin=19mm]{geometry}
\usepackage{amsmath,amssymb,mathtools}
\usepackage[scr=rsfs,cal=euler]{mathalfa}
\usepackage{xfrac}
\usepackage[unicode,hidelinks,bookmarks=false]{hyperref}
\newcommand{\ie}{i.e.\ }
\newcommand{\cf}{cf.\ }
\newcommand{\resp}{respectively\ }
\newcommand{\Subsection}{\subsection}
\providecommand{\nobreakdash}{\nobreak\hbox{-}}
\setlength{\emergencystretch}{2em}
\multlinegap0pt
\usepackage{amssymb}
\usepackage{mathtools}
\newtheorem{theorem}{Theorem}
\title{\vspace{-12mm}\textbf{Explicit Cartesian Coordinates for Three-Dimensional Clothoids}}
\author{Alexandru Ionu\textcommabelow{t}}
\date{Source: arXiv:2608.28636v1 (CC BY 4.0)}
\begin{document}
\maketitle

\noindent\emph{Source and licence.} \vspace{-12mm}\textbf{Explicit Cartesian Coordinates for Three-Dimensional Clothoids}. Version arXiv:2608.28636v1 (CC BY 4.0), distributed by arXiv under the Creative Commons Attribution 4.0 International licence, http://creativecommons.org/licenses/by/4.0/, as recorded in the arXiv licence metadata for this exact version and shown on the abstract page https://arxiv.org/abs/2608.28636v1. Full licence notice: \texttt{licenses/arxiv-2608.28636-CC-BY-4.0.txt}.\par\smallskip

\noindent\textit{Editorial scope.} Counted unit: the article's theorem theorem (source0.tex lines 352--379). The article's complete original proof of it is delivered with it (source0.tex lines 381--386, one \texttt{proof} environment of 6 lines). No mathematics is written, repaired or re-derived: the statement and the proof are the article's own lines, in the article's own notation, and no retained line is substituted or re-flowed.  No further source-local context is needed: every cross-reference of every retained line resolves inside the two delivered spans.  Nothing was omitted from any retained span, so the package carries no omission record.  No retained line contains a citation, so the wrapper carries no bibliography and no bibliography entry.  No retained line contains a figure environment, an image-inclusion command or drawing code, so no image is delivered and the package declares no asset dependency.  Editorial formatting only, in addition: an AMS \texttt{amsart} preamble that restates the article's own declarations and macros used by the retained lines, namely the environments \texttt{theorem} on separate counters, together with the standard packages those lines need.  The retained environments print in this document as Theorem 1 in order of appearance, whereas the article prints the same items with the numbering of its own full document.  The complete native source of the article as distributed in the arXiv e-print for this exact version is bundled unchanged as \texttt{sources/208799/source0.tex}.
\begin{theorem}[Explicit asymptotic lines]
Let the hypotheses of Theorem~1 hold, and set
\begin{equation}
 b_\pm=\rho^{-3/2}\left\{
 cC_\pm-cJ_\chi(\xi_0)
 -a\bigl(\xi_0\gamma(\xi_0)-\chi\mathbf d(\xi_0)\bigr)
 \right\}. \tag{42}
\end{equation}
Then, as \(s\to\pm\infty\),
\begin{align}
 r(s)={}&b_\pm+u_\pm\Bigg[
 \frac{a}{\rho}\left(s+\frac{ab+cd}{\rho^2}\right) \notag\\
 &\hspace{16mm}-\frac{c\Delta}{\rho^3}
 \log\left|\sqrt\rho\left(s+\frac{ab+cd}{\rho^2}\right)\right|
 \Bigg]+O(|s|^{-1}). \tag{43}
\end{align}
Consequently,
\begin{equation}
 \operatorname{dist}\bigl(r(s),\,b_\pm+\mathbb R u_\pm\bigr)
 =O(|s|^{-1}). \tag{44}
\end{equation}
The oriented angle \(\Theta\in[0,\pi]\) between the two asymptotic directions is
\begin{equation}
 \cos\Theta=2e^{-\pi\chi^2/2}-1,
 \qquad
 \Theta=2\arccos\!\left(e^{-\pi\chi^2/4}\right). \tag{45}
\end{equation}
\end{theorem}

\begin{proof}
Insert (33) and (39) into (16), use
\(\xi(s)=\sqrt\rho\bigl(s+(ab+cd)/\rho^2\bigr)\), and recall
\(\chi=-\Delta/\rho^{3/2}\). This gives (43) and the explicit basepoints (42). The logarithmic displacement is parallel to \(u_\pm\), so it does not affect the distance to the corresponding line, proving (44). Finally, (30) and the gamma-modulus identities give
\(u_+\cdot u_-=2e^{-\pi\chi^2/2}-1\), hence (45).
\end{proof}

\end{document}
